%% ============================================================
%%  CONCLUSION GÉNÉRALE ET PERSPECTIVES
%% ============================================================
\chapter*{Conclusion générale et perspectives}
\thispagestyle{plain}
\addcontentsline{toc}{chapter}{Conclusion générale et perspectives}
\markboth{Conclusion générale}{Conclusion générale}

\section*{Bilan du travail}

Ce projet de fin d'études avait pour objectif de remplacer
l'approche mono-modèle de génération de curricula par une
\textbf{famille d'agents spécialisés} reliés à une base de
connaissances pédagogique partagée à travers le Model Context
Protocol. Le présent rapport se concentre sur la
\textbf{conception logique} de cette famille d'agents et du
serveur qui les relie, indépendamment de toute pile technique
particulière. La stratégie d'orchestration qui articule la
coexistence de ces agents, ainsi que le dispositif d'observabilité
et l'analyse de sécurité associés, font l'objet du chapitre~3, qui
prolonge directement la présente conception.

Le travail s'organise autour de trois apports complémentaires.
Le premier est un \textbf{mécanisme d'indexation} des activités
pédagogiques, qui permet à un agent de consulter ce qui existe
déjà avant de produire du contenu nouveau. Le deuxième est un
\textbf{serveur \gls{mcp}} qui expose ces données et les
opérations associées sous la forme d'un catalogue d'outils
typés, indépendamment de la nature des clients qui les
consomment. Le troisième est un \textbf{catalogue d'agents} décrivant les
rôles retenus pour cette première version, aux responsabilités
distinctes et dont les privilèges sont restreints au strict
nécessaire par une liste blanche d'outils \gls{mcp}.

\section*{Apports du projet}

Les contributions principales de ce travail sont les suivantes.

\textbf{Une architecture orientée qualité.} Le déplacement du
travail d'un seul modèle vers une famille d'agents spécialisés a
été guidé par un objectif explicite, à savoir améliorer la
qualité du contenu produit, et non en accélérer la production.
La conception privilégie systématiquement l'ancrage des sorties
sur des données existantes, la séparation des privilèges et la
révision avant publication.

\textbf{Une indexation pensée pour la non-redondance.} Le
mécanisme d'indexation par plongement vectoriel n'est pas un
ajout périphérique mais un composant de premier plan : il
conditionne le comportement de l'agent générateur d'activités
avec \gls{mcp}, qui consulte la base avant toute production. La
représentation sémantique des activités et la recherche par
similarité s'enchaînent pour faire de la non-redondance une
propriété mesurable du système.

\textbf{Un protocole MCP en pratique.} L'usage du \gls{mcp}
permet de faire coexister, autour d'un même serveur, des agents
aux profils très différents, sans dupliquer la logique d'accès
aux données ni partager d'invite système. La discipline du contrat protocolaire
facilite l'évolution future : ajouter un agent supplémentaire,
par exemple l'assistant administrateur mentionné en perspective,
ne demandera aucune modification du serveur ni des agents
existants.

\textbf{Une lecture comparative des implémentations.}
Plutôt qu'une campagne de tests de conformité, le chapitre~2
propose une lecture comparative des formes d'agents livrées,
organisée autour de quatre axes : propriété de l'orchestration
(\gls{llm} contre code), posture humaine, granularité de
l'artefact inspectable, coût et latence. Cette lecture montre que
les variantes ne s'ordonnent pas sur une échelle linéaire de
qualité mais répondent à des questions d'intégration distinctes,
ce qui motive la stratégie d'orchestration développée au
chapitre~3.

\section*{Limites}

Plusieurs limites du présent travail méritent d'être signalées.

D'abord, l'\textbf{assistant administrateur}, mentionné en
perspective, n'est pas inclus dans la présente version. Sa
conception reprendra les principes des agents en lecture seule
(tuteur et assistant élève) mais étendus à des opérations de
supervision sur l'ensemble du curriculum.

Ensuite, le \textbf{choix définitif du modèle d'embedding} reste
ouvert au moment de la rédaction. Les compromis entre modèle
externe à \gls{api} et modèle local exécuté en interne ont été
documentés au chapitre 2, mais l'arbitrage final dépendra de
mesures de coût et de qualité à effectuer sur l'infrastructure
définitive.

Enfin, la lecture comparative du chapitre~2 s'appuie sur des
\textbf{observations d'ordre de grandeur} recueillies sur
l'instance déployée, et non sur un banc d'essai systématique. Une
évaluation quantitative à plus grande échelle — couverture des
objectifs pédagogiques, non-redondance mesurée, coûts par
variante — reste à mener et viendra étayer les constats
qualitatifs présentés ici.

\section*{Perspectives}

Plusieurs directions s'ouvrent pour la suite du travail.

La première concerne le \textbf{routage assisté entre variantes}.
Le chapitre~3 définit une stratégie d'orchestration où le choix de
la variante reste aujourd'hui explicite : l'enseignant choisit sa
surface. Une évolution naturelle consiste à \emph{suggérer}
automatiquement la variante adaptée à partir de l'intention
exprimée et du contexte de session, sans retirer à l'enseignant la
décision finale.

La deuxième concerne l'\textbf{ajout d'un agent supplémentaire},
l'assistant administrateur, dont la fonction de supervision
globale du curriculum complétera la famille existante. Sa
conception suit celle des agents en lecture seule (tuteur et
assistant élève) mais étendue à des opérations transversales :
détection d'activités obsolètes, statistiques de couverture des
objectifs, diagnostic d'incohérences entre unités.

La troisième concerne l'\textbf{évolution du modèle
d'embedding}. Le modèle initial sera choisi pour son coût et sa
simplicité, mais un modèle multilingue de plus haute dimension
pourrait être justifié si le volume d'activités croît
significativement, ou si la qualité du regroupement sémantique
devenait limitante en phase d'évaluation.

La quatrième concerne la \textbf{personnalisation des agents par
domaine}. Les fiches d'agent du chapitre 2 sont génériques. Une
spécialisation par discipline (mathématiques, sciences, langues)
permettrait de raffiner les rubriques d'évaluation qualitative
et les heuristiques d'ancrage propres à chaque discipline.

La cinquième et dernière concerne l'\textbf{approfondissement de
la sécurité}. Le chapitre~3 pose une analyse de sécurité
applicative (surface d'attaque \gls{llm}, contrôle d'accès,
gestion des secrets, traitement des données personnelles).
Les vecteurs propres aux systèmes agentiques — injection par
contenu indexé, empoisonnement de l'index, exfiltration par
recherche — appellent au-delà une étude offensive systématique
(tests d'intrusion ciblés).

\vspace{1cm}
En conclusion, ce projet montre qu'une famille d'agents
spécialisés, reliés par un contrat protocolaire standardisé et
disposant d'un mécanisme d'indexation commun, constitue une
réponse architecturale cohérente au problème de la génération
pédagogique. La valeur de cette architecture se mesure d'abord à
la qualité du contenu produit, à sa cohérence avec ce qui existe
déjà, et à sa résistance aux entrées adversariales. La validation
quantitative de cette affirmation, à l'échelle d'un banc d'essai
dédié, prolongera naturellement les constats qualitatifs établis
dans le présent document.
\newpage
